\chapter{Wstęp}
\label{ch:wstep}

\section{Wprowadzenie do problematyki kompresji dźwięku przy bardzo niskich przepływnościach}
\label{sec:wprowadzenie-kompresja-ultra-low}

Kompresja dźwięku polega na przedstawieniu sygnału akustycznego za pomocą mniejszej liczby bitów niż w reprezentacji pierwotnej. W przypadku audio można dążyć do wiernego odtworzenia przebiegu, zachowania jakości percepcyjnej albo przeniesienia informacji niezbędnej dla danego zastosowania. W kodowaniu mowy szczególne znaczenie ma ostatni z tych przypadków, ponieważ odbiorca ocenia przede wszystkim zrozumiałość wypowiedzi i naturalność głosu. Przy bardzo niskich przepływnościach problem zmienia charakter: sygnał mowy 16 kHz zapisany z rozdzielczością 16 bitów wymaga 256 kbps przed kompresją, a zejście poniżej 1 kbps oznacza ponad dwustukrotne ograniczenie ilości informacji. W takim reżimie kodek musi wybrać cechy najważniejsze dla rekonstrukcji użytecznej percepcyjnie.

Sygnał mowy jest podatny na taką selekcję, ponieważ posiada strukturę wynikającą z mechanizmu artykulacji. Istotne są między innymi zmiany widma w czasie, częstotliwość podstawowa, energia, formanty, fragmenty dźwięczne i bezdźwięczne oraz rytm wypowiedzi. Pominięcie zbyt wielu szczegółów prowadzi jednak do metalicznego brzmienia, utraty barwy głosu lub obniżenia zrozumiałości. W zakresie subkilobitowym celem nie jest więc klasyczna wierność sygnału, lecz kompromis między zrozumiałością, naturalnością i kosztem reprezentacji.

Klasyczne kodeki niskobitowe osiągają taki kompromis przez wykorzystanie modelu źródła mowy i przesyłanie niewielkiej liczby parametrów. Codec 2 oraz MELPe pokazują, że transmisja mowy poniżej lub blisko 1 kbps jest możliwa przy silnym uproszczeniu opisu sygnału \cite{Rowe2019Codec2,Rowe2021Codec2450,NATO2015STANAG4591}. Rozwój neuronowych metod kodowania wprowadził inne podejście: reprezentacja może być uczona bezpośrednio z danych, a koder i dekoder mogą wspólnie optymalizować wąskie gardło kompresyjne. SoundStream i EnCodec pokazały skuteczność autoenkoderów z dyskretną reprezentacją latentną i rezydualną kwantyzacją wektorową \cite{Zeghidour2021SoundStream,Defossez2022Encodec}. Wyzwanie polega jednak na przeniesieniu tych idei do przepływności poniżej 1 kbps, gdzie każdy dodatkowy bit wpływa na całkowity budżet transmisji.

\section{Motywacja do podjęcia tematu}
\label{sec:motywacja}

Motywacją do podjęcia tematu jest rosnące znaczenie systemów, w których sygnał mowy musi być przetwarzany, przesyłany lub archiwizowany przy ograniczonych zasobach. Dotyczy to zwłaszcza urządzeń brzegowych, takich jak terminale głosowe, czujniki akustyczne, urządzenia ubieralne oraz systemy pracujące w sieciach o małej przepustowości. W takich warunkach nie można zwykle poprawiać jakości przez proste zwiększanie bitrate'u. Budżet energetyczny, koszt transmisji, opóźnienie oraz pamięć stają się ograniczeniami równie ważnymi jak sama jakość rekonstrukcji.

Szczególnie interesujący jest zakres poniżej 1 kbps, ponieważ wymusza bardzo silną selekcję informacji przekazywanej przez kodek. W przypadku mowy najważniejsze jest zachowanie zrozumiałości oraz wybranych cech percepcyjnych, a nie pełnej wierności akustycznej. Klasyczne kodeki parametryczne, takie jak Codec 2 i MELPe, pokazują, że transmisja mowy w tym zakresie jest możliwa, lecz zwykle odbywa się kosztem naturalności brzmienia i odporności na zróżnicowane warunki akustyczne \cite{Rowe2019Codec2,Rowe2021Codec2450,NATO2015STANAG4591}. Uzasadnia to sprawdzenie, czy metody neuronowe mogą lepiej wykorzystać ograniczony budżet bitowy.

Drugim źródłem motywacji jest rozwój neuronowych kodeków audio. SoundStream i EnCodec pokazały, że autoenkoder z dyskretną reprezentacją latentną i rezydualną kwantyzacją wektorową może skutecznie kompresować sygnał audio \cite{Zeghidour2021SoundStream,Defossez2022Encodec}. Typowe warianty tych systemów są jednak projektowane dla przepływności większych niż zakres subkilobitowy i często wymagają większego budżetu obliczeniowego niż prostsze urządzenia brzegowe. Powstaje więc luka między lekkimi kodekami parametrycznymi a neuronowymi kodekami audio o wyższej jakości, lecz większym koszcie.

Temat pracy wynika z potrzeby zbadania tej luki w sposób praktyczny. Istotna jest nie tylko mała nominalna przepływność, lecz także zgodność architektury, kwantyzacji i kosztu inferencji z ograniczeniami przetwarzania brzegowego. Kierunek ten ma również znaczenie dla prywatności, ponieważ lokalne kodowanie mowy może ograniczyć potrzebę przesyłania pełnego sygnału audio do infrastruktury zewnętrznej. W pracy przyjęto więc perspektywę eksperymentalną: oceniono, czy zwarta reprezentacja neuronowa może zachować wystarczającą informację mowy poniżej 1 kbps i pozostać możliwa do uruchomienia w lekkim środowisku wykonawczym.

\section{Cel i zakres pracy}
\label{sec:cel-zakres}

Celem pracy jest opracowanie, implementacja i ocena neuronowego kodeka mowy pracującego przy ultra-niskiej przepływności bitowej, z uwzględnieniem ograniczeń typowych dla przetwarzania brzegowego. Przez ultra-niską przepływność rozumiany jest w pracy zakres poniżej 1 kbps, w którym klasyczne zachowanie pełnej informacji akustycznej staje się bardzo trudne, a istotne jest przede wszystkim przeniesienie informacji potrzebnej do zrozumienia wypowiedzi oraz uzyskania możliwie naturalnej rekonstrukcji sygnału mowy. Proponowane rozwiązanie ma charakter eksperymentalny: jego celem nie jest zastąpienie dojrzałych standardów telekomunikacyjnych, lecz sprawdzenie, w jakim stopniu architektura oparta na autoenkoderze, reprezentacji latentnej i rezydualnej kwantyzacji wektorowej może zostać dostosowana do pracy w reżimie subkilobitowym.

\subsection{Cele szczegółowe}
\label{subsec:cele-szczegolowe}

Cele szczegółowe pracy można ująć w następujących punktach:

\begin{itemize}
    \item \textbf{Analiza stanu wiedzy w zakresie ultra-niskobitowego kodowania mowy.} Celem było określenie, jakie rozwiązania klasyczne, hybrydowe i neuronowe są stosowane przy bardzo niskich przepływnościach oraz jakie ograniczenia pojawiają się przy zejściu poniżej 1 kbps.
    \item \textbf{Określenie wymagań projektowych kodeka.} W pracy przyjęto wymagania dotyczące docelowej przepływności bitowej, jakości rekonstrukcji, dopuszczalnego opóźnienia, stabilności kwantyzacji oraz możliwości uruchomienia modelu poza środowiskiem treningowym.
    \item \textbf{Zaprojektowanie architektury neuronowego kodeka mowy.} Obejmowało to przygotowanie koncepcji kodera, reprezentacji latentnej, modułu rezydualnej kwantyzacji wektorowej oraz dekodera rekonstruującego przebieg czasowy sygnału.
    \item \textbf{Przygotowanie środowiska eksperymentalnego.} Istotnym celem było opracowanie pipeline'u danych, konfiguracji treningu, mechanizmów zapisywania checkpointów, logowania metryk oraz generowania przykładów zrekonstruowanego dźwięku.
    \item \textbf{Przeprowadzenie eksperymentów i ocena jakości rekonstrukcji.} W ramach pracy porównano warianty modelu, analizowano wpływ parametrów architektury i kwantyzacji na jakość odtwarzanej mowy oraz oceniano artefakty obecne w rekonstruowanym sygnale.
    \item \textbf{Ocena przydatności rozwiązania w przetwarzaniu brzegowym.} Sprawdzono koszt inferencji, wykorzystanie pamięci oraz możliwość przygotowania modelu do uruchomienia w lekkim środowisku wykonawczym.
\end{itemize}

\subsection{Zakres pracy}
\label{subsec:zakres-pracy}

Zakres pracy obejmuje następujące obszary:

\begin{itemize}
    \item \textbf{Zakres teoretyczny.} W pracy omówiono podstawy kodowania mowy i dźwięku przy bardzo niskich przepływnościach, charakterystykę sygnału mowy, reprezentację cyfrową audio, autoenkodery, konwolucyjne przetwarzanie sygnału, reprezentacje latentne, kwantyzację wektorową oraz metryki oceny jakości rekonstrukcji.
    \item \textbf{Zakres implementacyjny.} Przygotowano prototyp kodeka składającego się z konwolucyjnego kodera, modułu rezydualnej kwantyzacji wektorowej oraz dekodera rekonstruującego przebieg czasowy. Uwzględniono również narzędzia pomocnicze potrzebne do treningu, konfiguracji eksperymentów i analizy wyników.
    \item \textbf{Zakres eksperymentalny.} Badano zależności między przepływnością bitową, doborem słowników RVQ, parametrami architektury oraz jakością odtwarzanej mowy. Ocena wyników została oparta na metrykach obiektywnych, analizie spektrogramów oraz odsłuchowej identyfikacji artefaktów.
    \item \textbf{Zakres wdrożeniowy.} Praca obejmuje analizę kosztu inferencji, wykorzystania pamięci i możliwości przygotowania modelu do uruchomienia na platformie brzegowej, w tym w środowisku wykonawczym przeznaczonym dla lekkich modeli.
    \item \textbf{Ograniczenia zakresu.} Poza zakresem pracy pozostaje budowa pełnego standardu transmisyjnego, implementacja zaawansowanego kodowania entropijnego jako produkcyjnego elementu strumienia bitowego oraz przeprowadzenie formalnych testów odsłuchowych z dużą grupą uczestników.
\end{itemize}

Tak określony zakres pozwala skoncentrować się na najważniejszym problemie pracy: ocenie, czy bardzo zwarta reprezentacja neuronowa może zachować użyteczną informację mowy przy przepływności poniżej 1 kbps.

Struktura pracy jest następująca. W rozdziale 2 przedstawiono podstawy teoretyczne dotyczące kodowania mowy, neuronowej kompresji sygnału oraz oceny jakości rekonstrukcji. Rozdział 3 opisuje założenia projektowe, architekturę proponowanego kodeka, proces kompresji i rekonstrukcji oraz wykorzystane technologie. W rozdziale 4 omówiono zbiór danych, konfigurację eksperymentów, wpływ parametrów modelu na jakość rekonstrukcji oraz wyniki ewaluacji. Rozdział 5 zawiera wnioski z przeprowadzonych badań, ocenę realizacji celu pracy, omówienie ograniczeń technicznych oraz możliwe kierunki dalszego rozwoju.

\section{Podział prac}
\label{sec:podzial-prac}

Praca miała charakter zespołowy, dlatego podział zadań należy rozumieć jako podział odpowiedzialności merytorycznej, a nie jako prosty podział według historii zmian w repozytorium. Część decyzji projektowych, implementacyjnych i eksperymentalnych była podejmowana wspólnie podczas pracy przy jednym stanowisku, zwłaszcza przy strojeniu architektury kodeka oraz interpretacji wyników odsłuchowych i metrycznych.

Do zadań wspólnych należał przegląd literatury dotyczącej neuronowych kodeków mowy, w szczególności metod opartych na autoenkoderach, kwantyzacji wektorowej i uczeniu reprezentacji dyskretnych. Wspólnie opracowano koncepcję kodeka oraz wymagania projektowe obejmujące przepływność bitową, jakość rekonstrukcji, opóźnienie i warianty działania systemu. Na tej podstawie przygotowano ogólną architekturę, obejmującą koder sygnału, reprezentację latentną, moduł kwantyzacji oraz dekoder fali dźwiękowej. Wspólnie określono również protokół ewaluacji i sformułowano końcowe wnioski.

Jakub Aszyk odpowiadał za koder wejściowy, przygotowanie danych oraz wdrożenie systemu na platformie brzegowej. W ramach implementacji opracował konwolucyjny tor kodera, który przekształcał surową falę mowy do krótszej sekwencji wektorów latentnych przekazywanych do kwantyzacji. Obejmowało to dobór redukcji czasowej, szerokości kanałów, wymiaru reprezentacji latentnej oraz sposobu przygotowania wyjścia kodera pod moduł RVQ. Część ta bezpośrednio wpływała na nominalną przepływność bitową, pojemność reprezentacji i koszt dalszego dekodowania. Jakub przygotował także pipeline danych, konfiguracje eksperymentów, logowanie przebiegów treningu oraz porządkowanie artefaktów, takich jak checkpointy, spektrogramy i przykłady rekonstrukcji. W części brzegowej odpowiadał za analizę czasu inferencji, wykorzystania pamięci, stabilności modelu, przygotowanie wersji eksportowalnej, konwersję do TensorFlow Lite oraz testy zgodności z implementacją referencyjną.

Kosma Gąsiorowski odpowiadał za część modelową związaną z rekonstrukcją i kwantyzacją oraz za zasadniczą część badań jakościowych. Rozwijał moduł rezydualnej kwantyzacji wektorowej, dekoder rekonstruujący przebieg czasowy oraz warianty poprawiające jakość odtwarzania, w tym bloki temporalne w przestrzeni latentnej, elementy samouwagi, dodatkowe głowice dekodera, aktywacje okresowe i moduł rekonstrukcji wysokich częstotliwości. Przygotował również funkcje kosztu używane podczas uczenia, obejmujące straty czasowe, wieloskalowe straty STFT, straty melowe, składniki stabilizujące kwantyzację oraz kary związane z użyciem słowników kodowych. W części eksperymentalnej odpowiadał za strojenie hiperparametrów, porównania wariantów architektury, analizę ablacyjną, ocenę jakości odtworzonego dźwięku i interpretację błędów na podstawie metryk, spektrogramów oraz testów porównawczych.
